\chapter{Kurva}\label{ch:b2:curves}

Kini geometrinya berbelok menjadi \omterm{def:b2:diffcalc:differential}{diferensial}. Sebuah kurva adalah titik yang bergerak menembus
ruang; kalkulus memberi kita kecepatan dan percepatannya, sedangkan geometri
bertanya apa yang hakiki --- yakni yang tak bergantung pada secepat apa kita menyusuri
lintasannya. Jawabannya adalah \emph{\omterm{def:b2:curves:length}{panjang busur}}, yang mengukur
lintasannya sendiri, dan \emph{\omterm{thm:b2:curves:frenet2d}{kelengkungan}}, yang mengukur seberapa ia
membengkok. Pada dimensi $3$ ada invarian kedua, yakni \emph{\omterm{thm:b2:curves:frenet3d}{torsi}},
yang mengukur seberapa kurvanya terpuntir keluar dari bidangnya. Adapun alat tata bukunya
untuk semua ini adalah \emph{kerangka Frenet} yang bergerak.

\section{Busur terparameter}

\begin{definition}[Busur terparameter]\label{def:b2:curves:arc}
Sebuah \emph{busur terparameter}\index{busur terparameter} berkelas
$\mathcal{C}^k$ (dengan $k \geq 1$) adalah
pemetaan $\gamma \colon I \to \R^n$ berkelas $\mathcal{C}^k$ pada
interval $I$. Sebuah titik $\gamma(t)$ disebut \emph{regular} bila
$\gamma'(t) \neq 0$, dan busurnya \emph{regular} bila semua
titiknya demikian. Adapun garis lewat $\gamma(t)$ yang diarahkan $\gamma'(t)$
disebut \emph{garis singgung}\index{garis singgung} di titik regularnya.
\end{definition}

\begin{definition}[Penggantian parameter]\label{def:b2:curves:reparam}
Sebuah \emph{penggantian parameter}\index{penggantian parameter} berkelas
$\mathcal{C}^k$ adalah
difeomorfisma $\mathcal{C}^k$ $\theta \colon J \to I$ antar
interval (dengan $\theta' \neq 0$ di mana-mana). Busur $\gamma$ dan
$\gamma \circ \theta$ disebut \emph{setara}; sedangkan sebuah
\emph{busur geometris}\index{busur geometris} (atau kurva) adalah sebuah kelas
kesetaraan. Adapun gagasan yang
invarian terhadap penggantian parameter --- yakni lintasannya, \omterm{def:b2:curves:arc}{garis singgungnya},
\omterm{def:b2:curves:length}{panjang busurnya}, \omterm{thm:b2:curves:frenet2d}{kelengkungannya} --- disebut \emph{geometris}.
\end{definition}

\begin{remark}
Lintasannya sendiri tak menentukan busur geometrisnya: sebab
parameterisasi $t \mapsto (\cos t, \sin t)$ pada $[0, 2\pi]$ dan pada
$[0, 4\pi]$ punya peta yang sama tetapi menempuh lingkarannya sekali dan
dua kali. Jadi sebuah \omterm{def:b2:curves:reparam}{busur geometris} mengingat multiplisitas dan orientasi
penempuhannya, bukan kelajuannya.
\end{remark}

\begin{example}[Kelajuan tak mengubah apa pun yang
geometris]\label{ex:b2:curves:speedinvariance}
Parameterkan lingkaran satuannya lewat
\[
\gamma(t) = (\cos t^2,\ \sin t^2),
\qquad t \in \intcc0{\sqrt{2\pi}} .
\]
Kelajuannya $\norm{\gamma'(t)} = 2t$ tumbuh linear, namun
\[
L = \int_0^{\sqrt{2\pi}}2t\,\dd t = 2\pi ,
\]
yakni \omterm{def:b2:curves:length}{panjang} yang sama seperti pada kelajuan tetap --- sebagaimana
dijanjikan \cref{thm:b2:curves:lengthinv}, lewat \omterm{def:b2:curves:reparam}{penggantian
parameter} $t \mapsto t^2$. \omterm{def:b2:curves:arc}{Garis singgungnya}, \omterm{thm:b2:curves:frenet2d}{kelengkungan} yang
dihitung dari \cref{prop:b2:curves:kappaformula}, dan setiap
besaran geometris lainnya juga sepakat; hanya $t = 0$
yang pantas dilirik, tempat $\gamma'(0) = 0$ membuat \emph{parameterisasi
ini} tak regular walaupun lintasannya berupa
lingkaran yang sempurna. Jadi pernyataan geometris kurang toleran terhadap
parameterisasi yang buruk: parameterkan ulang dulu, simpulkan
kemudian.
\end{example}

\begin{example}\label{ex:b2:curves:cusp}
Busur $\gamma(t) = (t^2, t^3)$ bersifat $\mathcal{C}^\infty$ tetapi tak
regular: sebab $\gamma'(0) = (0, 0)$. Lintasannya, yakni parabola
semikubik $y^2 = x^3$, punya \emph{titik runcing} di titik asalnya: jadi kemulusan
parameterisasinya tak mencegah kesingularan geometris
di tempat kecepatannya lenyap. Itulah sebabnya hipotesis keregularan
$\gamma' \neq 0$ bukan sekadar hiasan.
\end{example}

\section{Panjang busur}

\begin{definition}[Panjang busur]\label{def:b2:curves:length}
Misalkan $\gamma \colon [a, b] \to \R^n$ sebuah busur $\mathcal{C}^1$. Adapun
\emph{panjang}\index{panjang busur} busurnya adalah
\[
L(\gamma) = \int_a^b \norm{\gamma'(t)}\, \dd t ,
\]
dengan $\norm{\cdot}$ menyatakan \omterm{def:b2:nvs:norm}{norma} Euklides. Sedangkan \emph{fungsi panjang
busur}\index{fungsi panjang busur} yang berpangkal di $t_0$ adalah $s(t) = \int_{t_0}^t
\norm{\gamma'(u)}\,\dd u$.
\end{definition}

\begin{theorem}[Panjangnya geometris; pencirian
poligonal]\label{thm:b2:curves:lengthinv}
\leavevmode
\begin{enumerate}
\item Bila $\theta \colon [c, d] \to [a, b]$ merupakan \omterm{def:b2:curves:reparam}{penggantian
parameter} $\mathcal{C}^1$, maka $L(\gamma \circ \theta) = L(\gamma)$.
\item Adapun $L(\gamma)$ merupakan supremum \omterm{def:b2:curves:length}{panjang} poligon yang
terlukis di dalamnya:
\[
L(\gamma)
= \sup\Bigl\{\, \sum_{i=1}^{m} \norm{\gamma(t_i) -
\gamma(t_{i-1})} \;:\; a = t_0 < t_1 < \dots < t_m = b \,\Bigr\}.
\]
\end{enumerate}
\end{theorem}

\begin{proof}
\emph{1.} Lewat penggantian peubah $t = \theta(u)$ (jilid Tahun
ke-1, yang sah karena $\theta$ bersifat $\mathcal{C}^1$ monoton),
\[
\int_c^d \norm{(\gamma\circ\theta)'(u)}\,\dd u
= \int_c^d \norm{\gamma'(\theta(u))}\,\abs{\theta'(u)}\,\dd u
= \int_a^b \norm{\gamma'(t)}\,\dd t ,
\]
tempat kita memakai $(\gamma\circ\theta)' = \theta'\cdot
(\gamma'\circ\theta)$ dan, bila $\theta$ turun, tanda
$\theta'$ terserap oleh pembalikan batasnya.

\emph{2.} Untuk sembarang subbagiannya, $\gamma(t_i) - \gamma(t_{i-1}) =
\int_{t_{i-1}}^{t_i}\gamma'(t)\,\dd t$, jadi menurut ketaksamaan
segitiga bagi integral berlaku $\norm{\gamma(t_i) - \gamma(t_{i-1})}
\leq \int_{t_{i-1}}^{t_i}\norm{\gamma'}$: jadi setiap poligonnya lebih pendek
daripada $L(\gamma)$, sehingga $\sup \leq L(\gamma)$.

Untuk ketaksamaan sebaliknya, misalkan $\varepsilon > 0$. Karena $\gamma'$
\omterm{def:b2:metric:continuity}{kontinu} pada $[a,b]$ yang \omterm{def:b2:metric:compact}{kompak}, ia kontinu seragam:
jadi ada $\delta > 0$ dengan $\norm{\gamma'(t) - \gamma'(u)} \leq
\varepsilon$ setiap kali $\abs{t - u} \leq \delta$. Ambillah sebuah subbagian
berlangkah $\leq \delta$. Pada tiap kepingnya, untuk $t \in [t_{i-1}, t_i]$,
\[
\gamma(t_i) - \gamma(t_{i-1})
= \int_{t_{i-1}}^{t_i}\gamma'(t)\,\dd t
= (t_i - t_{i-1})\,\gamma'(t_{i-1}) + R_i,
\qquad
\norm{R_i} \leq \varepsilon\,(t_i - t_{i-1}),
\]
karena $R_i = \int_{t_{i-1}}^{t_i}(\gamma'(t) -
\gamma'(t_{i-1}))\,\dd t$. Karena itu
\[
\norm{\gamma(t_i) - \gamma(t_{i-1})}
\geq (t_i - t_{i-1})\norm{\gamma'(t_{i-1})}
- \varepsilon (t_i - t_{i-1}) .
\]
Setelah dijumlahkan, lalu membandingkan $\sum (t_i -
t_{i-1})\norm{\gamma'(t_{i-1})}$ dengan $\int_a^b\norm{\gamma'}$ (yakni
jumlah Riemann bagi fungsi \omterm{def:b2:metric:continuity}{kontinu} $\norm{\gamma'}$, yang berselisih paling
$\varepsilon(b - a)$ dari integralnya untuk $\delta$ yang cukup kecil berkat
kekontinuan \omterm{def:b2:funcseq:def}{seragamnya} lagi), kita memperoleh poligon berpanjang $\geq
L(\gamma) - 2\varepsilon(b - a)$. Lalu melewatkan $\varepsilon \to 0$
membuktikan klaimnya.
\end{proof}

\begin{example}[Archimedes dan poligon yang
terlukis]\label{ex:b2:curves:archimedes}
Untuk lingkaran satuannya, segi-$n$ beraturan yang terlukis di dalamnya berpanjang
$L_n = 2n\sin\frac\pi n$, dan uraian $\sin x = x -
\frac{x^3}6 + O(x^5)$ memberi
\[
L_n = 2\pi - \frac{\pi^3}{3n^2} + O\Bigl(\frac1{n^4}\Bigr) :
\]
jadi \omterm{def:b2:curves:length}{panjang} poligonal pada \cref{thm:b2:curves:lengthinv}
konvergen secara \emph{kuadratik}. Secara numerik: $L_6 = 6$
(yakni segi enamnya, yang memberi $\pi > 3$ yang kasar), sedangkan $L_{96} =
192\sin\frac{\pi}{96} \approx 6.28206$ terhadap $2\pi \approx
6.28319$ --- dan galatnya $0.00113$ sepakat dengan ramalan
$\pi^3/(3\cdot96^2) \approx 0.00112$. Itulah sebabnya Archimedes,
dengan melipatduakan segi enamnya lima kali menjadi $96$ sisi, dapat mengapit
$\pi$ sampai tiga digit dengan tangan: sebab tiap pelipatduaannya membagi
galatnya dengan empat. Jadi supremum pada pencirian
poligonalnya bukan sekadar tercapai pada limitnya; ia
tercapai secara \emph{cepat}, karena kurva yang mulus memisah dari
tali busurnya hanya pada orde kedua.
\end{example}

\begin{theorem}[Parameterisasi panjang busur]\label{thm:b2:curves:arclength}
Misalkan $\gamma \colon I \to \R^n$ busur $\mathcal{C}^k$ yang regular
(dengan $k \geq 1$). Adapun \omterm{def:b2:curves:length}{fungsi panjang busurnya} $s$ merupakan difeomorfisma
$\mathcal{C}^k$ dari $I$ ke sebuah interval $J$, dan $\tilde\gamma =
\gamma \circ s^{-1}$ memenuhi $\norm{\tilde\gamma'} = 1$
di mana-mana. Hingga penggeseran parameternya dan orientasinya,
parameterisasi \emph{\omterm{def:b2:curves:length}{panjang busur}} (atau \emph{berkelajuan satuan}) ini
bersifat tunggal.
\end{theorem}

\begin{proof}
Berlaku $s'(t) = \norm{\gamma'(t)} > 0$ dan $s'$ bersifat $\mathcal{C}^{k-1}$
(yakni komposisi pemetaan $\mathcal{C}^{k-1}$ $\gamma'$ dengan \omterm{def:b2:nvs:norm}{normanya},
yang mulus jauh dari $0$), jadi $s$ bersifat $\mathcal{C}^k$, naik
tegas, bijeksi ke $J = s(I)$, dan balikannya bersifat
$\mathcal{C}^k$ menurut teorema fungsi balikan pada dimensi $1$
(jilid Tahun ke-1). Maka
\[
\tilde\gamma'(\sigma)
= \frac{\gamma'(t)}{s'(t)}
= \frac{\gamma'(t)}{\norm{\gamma'(t)}},
\qquad t = s^{-1}(\sigma),
\]
yakni sebuah vektor satuan. Bila $\hat\gamma = \gamma\circ\theta$ merupakan
parameterisasi berkelajuan satuan yang lain, maka $\abs{\theta'} = 1$, jadi $\theta'
= \pm 1$ yang tetap (berkat kekontinuannya), yakni $\theta(u) = \pm u + c$.
\end{proof}

\begin{remark}
\omterm{def:b2:curves:length}{Panjang busur} adalah parameter yang memisahkan geometri dari
dinamika. Sebuah lintasan dapat ditempuh dengan profil kelajuan
apa pun --- yakni fisika geraknya --- tetapi setiap pertanyaan yang
invarian terhadap parameterisasi (yakni bentuknya, pembengkokannya,
oskulasinya) punya jam kanonik, yaitu jarak yang ditempuh.
Itulah sebabnya semua rumus \omterm{thm:b2:curves:frenet2d}{kelengkungan} di bawah \emph{didefinisikan}
pada kelajuan satuan lalu \emph{diterjemahkan} ke parameterisasi
sembarang lewat aturan rantai: sebab faktor terjemahannya
adalah pangkat $v = s'$, dan melacaknya secara benar adalah
seluruh isi \cref{prop:b2:curves:kappaformula}.
\end{remark}

\begin{example}[Lingkaran dan heliks]\label{ex:b2:curves:helix}
Untuk lingkaran $\gamma(t) = (R\cos t, R\sin t)$, berlaku $\norm{\gamma'} =
R$, jadi $s = Rt$ dan \omterm{def:b2:curves:length}{panjang} satu putaran penuhnya adalah $2\pi R$. Sedangkan untuk
heliks $\gamma(t) = (a\cos t,\ a\sin t,\ bt)$ dengan $a > 0$,
$\norm{\gamma'(t)} = \sqrt{a^2 + b^2}$ bersifat tetap: jadi heliksnya
ditempuh pada kelajuan tetap, dan $s = t\sqrt{a^2 + b^2}$.
\end{example}

\begin{example}[Panjang busur dalam koordinat
kutub]\label{ex:b2:curves:polarlength}
Sebuah kurva kutub $r = r(\theta)$ adalah busur $\gamma(\theta) =
(r\cos\theta,\ r\sin\theta)$, dengan
\[
\gamma'(\theta) = (r'\cos\theta - r\sin\theta,\
r'\sin\theta + r\cos\theta),
\qquad
\norm{\gamma'(\theta)}^2 = r'^2 + r^2
\]
(sebab suku silangnya saling meniadakan): jadi unsur \omterm{def:b2:curves:length}{panjang} kutubnya adalah
$\sqrt{r^2 + r'^2}\,\dd\theta$. Untuk kardioid $r = 1 +
\cos\theta$:
\[
r^2 + r'^2 = (1 + \cos\theta)^2 + \sin^2\theta = 2 +
2\cos\theta = 4\cos^2\tfrac\theta2,
\]
lalu pada $\intcc{-\pi}{\pi}$ kosinus setengah sudutnya bersifat
taknegatif, jadi
\[
L = \int_{-\pi}^{\pi}2\cos\tfrac\theta2\,\dd\theta
= \Bigl[4\sin\tfrac\theta2\Bigr]_{-\pi}^{\pi}
= 4 - (-4) = 8 :
\]
seperti lengkung sikloid pada \cref{exo:b2:curves:1}, sebuah kurva
yang \omterm{def:b2:structures:generated}{dibangun} dari lingkaran punya \omterm{def:b2:curves:length}{panjang} yang rasional, tanpa $\pi$
di mana pun. Adapun pemfaktoran setengah sudutnya adalah siasat baku
bagi \omterm{def:b2:curves:length}{panjang} kurva yang terbangkitkan lingkaran; sedangkan ketika ia gagal (yakni pada
elipsnya), \omterm{def:b2:curves:length}{panjangnya} menjadi fungsi yang sungguh baru --- yakni sebuah
integral eliptik, di luar bentuk tertutup yang dasar.
\end{example}

\section{Kelengkungan pada bidang}

Di sepanjang bagian ini, busurnya bersifat $\mathcal{C}^2$ dan regular pada
bidang Euklides yang berorientasi. Kita memparameterkannya lewat \omterm{def:b2:curves:length}{panjang busur} lalu
menulis $T(s) = \tilde\gamma'(s)$ bagi singgung satuannya, dan $N(s)$
bagi vektor satuan yang langsung ortogonal terhadap $T(s)$ (yakni perputaran
$T$ sebesar $+\pi/2$).

\begin{theorem}[Rumus Frenet bidang]\label{thm:b2:curves:frenet2d}
Misalkan $\tilde\gamma$ busur $\mathcal{C}^2$ berkelajuan satuan pada
bidang yang berorientasi. Maka ada fungsi \omterm{def:b2:metric:continuity}{kontinu} $\kappa$, yakni
\emph{kelengkungan}\index{kelengkungan} (aljabarnya), sedemikian sehingga
\[
T'(s) = \kappa(s)\, N(s),
\qquad
N'(s) = -\kappa(s)\, T(s).
\]
\end{theorem}

\begin{proof}
Karena $\norm{T(s)}^2 = 1$ untuk setiap $s$, menurunkan hasil kali
skalarnya memberi $2\langle T'(s), T(s)\rangle = 0$: jadi $T'(s)$ ortogonal terhadap
$T(s)$, sehingga segaris dengan $N(s)$ (berkat dimensi $2$); jadi tulislah $T'(s)
= \kappa(s) N(s)$ dengan $\kappa(s) = \langle T'(s), N(s)\rangle$,
yang \omterm{def:b2:metric:continuity}{kontinu}. Demikian pula $N' \perp N$, jadi $N' = \lambda T$; lalu
menurunkan $\langle T, N\rangle = 0$ memberi $\langle T', N\rangle +
\langle T, N'\rangle = \kappa + \lambda = 0$.
\end{proof}

\begin{definition}\label{def:b2:curves:curvature}
Ketika $\kappa(s) \neq 0$, \emph{jari-jari
kelengkungan}\index{jari-jari kelengkungan} adalah $R(s) =
1/\abs{\kappa(s)}$ dan \emph{pusat kelengkungan}\index{pusat
kelengkungan} adalah
$\tilde\gamma(s) + \frac{1}{\kappa(s)} N(s)$; sedangkan lingkaran dengan pusat
itu dan berjari-jari $R(s)$ disebut \emph{lingkaran oskulasi}\index{lingkaran oskulasi}, yakni
hampiran lingkaran terbaik bagi kurvanya di $\tilde\gamma(s)$.
\end{definition}

\begin{example}[Lingkaran oskulasi bagi
eksponensialnya]\label{ex:b2:curves:oscexp}
Untuk $y = \eu^x$ di titik $(0, 1)$: berlaku $f'(0) = f''(0) = 1$,
jadi menurut rumus grafik di bawah,
\[
\kappa(0) = \frac{1}{(1 + 1)^{3/2}} = \frac1{2\sqrt2},
\qquad R = 2\sqrt2 .
\]
Adapun singgung satuannya $T = \frac{(1, 1)}{\sqrt2}$, normal
langsungnya $N = \frac{(-1, 1)}{\sqrt2}$, dan \omterm{def:b2:curves:curvature}{pusat
kelengkungannya} adalah
\[
(0, 1) + 2\sqrt2\cdot\frac{(-1, 1)}{\sqrt2} = (-2,\ 3) :
\]
jadi \omterm{def:b2:curves:curvature}{lingkaran oskulasinya} berpersamaan $(x + 2)^2 + (y - 3)^2 =
8$. Sebagai pemeriksaan klaim ``hampiran lingkaran terbaik'' itu:
menyelesaikan persamaan lingkarannya untuk $y$ di dekat $(0,1)$ lalu
menguraikannya memberi $y = 1 + x + \frac{x^2}2 + O(x^3)$ ---
persis uraian Taylor orde dua bagi $\eu^x$. Jadi
\omterm{def:b2:curves:curvature}{lingkaran oskulasinya} mencocokkan nilai, lereng \emph{dan} turunan
kedua; sedangkan lingkaran singgung biasa hanya akan mencocokkan
dua yang pertama.
\end{example}

\begin{example}[Evolut sebuah lingkaran adalah
pusatnya]\label{ex:b2:curves:circleevolute}
Untuk lingkaran berjari-jari $R$ yang ditempuh berlawanan arah jarum jam,
berlaku $\kappa = 1/R$ dan $N$ menunjuk ke pusatnya, jadi
\omterm{def:b2:curves:curvature}{pusat kelengkungannya} $\tilde\gamma + \frac1\kappa N$ adalah
pusat lingkarannya, untuk setiap $s$: jadi \omterm{def:b2:curves:curvature}{lingkaran oskulasi}
sebuah lingkaran adalah lingkarannya sendiri, dan tempat kedudukan \omterm{def:b2:curves:curvature}{pusat
kelengkungannya} runtuh menjadi satu titik. Kasus merosot ini
menakar \cref{exo:b2:curves:6}: sebab di sana kecepatan evolutnya
adalah $-\frac{\kappa'}{\kappa^2}N$, yang lenyap
secara identik persis ketika $\kappa$ bersifat tetap.
\end{example}

\begin{proposition}[Kelengkungan dalam parameterisasi
sembarang]\label{prop:b2:curves:kappaformula}
Untuk busur bidang $\mathcal{C}^2$ yang regular $\gamma(t) = (x(t), y(t))$,
\[
\kappa(t)
= \frac{x'(t)\,y''(t) - y'(t)\,x''(t)}
       {\bigl(x'(t)^2 + y'(t)^2\bigr)^{3/2}} ,
\]
khususnya $\kappa = \dfrac{y''}{(1 + y'^2)^{3/2}}$ bagi sebuah grafik
$y = f(x)$.
\end{proposition}

\begin{proof}
Tulislah $v(t) = \norm{\gamma'(t)} = s'(t)$, jadi $\gamma' = vT$
(dengan mengomposisikan data berkelajuan satuannya dengan $s$). Setelah diturunkan,
\[
\gamma'' = v'T + v\,T'\cdot s' = v'T + v^2\kappa N .
\]
Kini ambillah \omterm{def:b2:linalg:det}{determinan} (dalam basis kanonik yang berorientasi) atas
$(\gamma', \gamma'')$: sebab $\det(T, T) = 0$ dan $\det(T, N) = 1$,
\[
\det(\gamma', \gamma'') = \det(vT,\ v'T + v^2\kappa N)
= v^3\kappa .
\]
Ruas kirinya adalah $x'y'' - y'x''$, dan $v^3 = (x'^2 +
y'^2)^{3/2}$. Adapun kasus grafiknya adalah parameterisasi $t \mapsto (t,
f(t))$.
\end{proof}

\begin{example}[Garis, lingkaran, parabola]\label{ex:b2:curves:kappaexamples}
Sebuah garis punya $\kappa = 0$ (dan sebaliknya: sebab $T' = 0$ berarti $T$
tetap, jadi $\tilde\gamma(s) = \tilde\gamma(0) + sT$, yakni sebuah garis). Adapun
lingkaran berjari-jari $R$ yang ditempuh berlawanan arah jarum jam punya $\kappa = 1/R$:
sebab dengan $\gamma(t) = (R\cos t, R\sin t)$, rumusnya memberi $\kappa =
R^2/R^3$. Sedangkan untuk parabola $y = x^2/2$: $\kappa(x) = 1/(1 +
x^2)^{3/2}$, yang maksimum di puncaknya --- jadi parabolanya paling tajam
membengkok di tempat ia berbalik.
\end{example}

\begin{remark}[Jebakan yang sering muncul seputar kelengkungan]
(i) \omterm{thm:b2:curves:frenet2d}{Kelengkungan} aljabar sebuah busur bidang berganti tanda ketika
orientasi busurnya atau bidangnya dibalik: jadi hanya
$\abs\kappa$ dan $R = 1/\abs\kappa$ yang murni geometris. Sebuah
lingkaran yang ditempuh searah jarum jam punya $\kappa = -1/R$. (ii) Adapun
rumus grafiknya $\kappa = f''/(1 + f'^2)^{3/2}$ diam-diam
memilih parameterisasi lewat $x$; jadi menerapkannya pada kurva
yang bukan grafik di dekat titiknya (yakni bersinggung tegak) adalah
kekeliruan yang klasik. (iii) Di titik tempat $\gamma' = 0$
tak ada yang terdefinisi --- baik $T$ maupun $\kappa$ --- dan
lintasannya boleh jadi sungguh patah
(\cref{ex:b2:curves:cusp}); jadi periksalah selalu keregularannya sebelum
menurunkan singgung satuannya. (iv) Di ruang, $\kappa =
\norm{T'} \geq 0$ menurut konvensinya: jadi tak ada tanda yang bisa
salah, tetapi juga tak ada tanda yang bisa dieksploitasi --- sebab
informasi bertipe belokan pindah ke \omterm{thm:b2:curves:frenet3d}{torsinya}. (v) Akhirnya, $\kappa$ adalah
turunan \emph{terhadap \omterm{def:b2:curves:length}{panjang busurnya}}: jadi untuk
parameterisasi yang tak berkelajuan satuan, melupakan faktor $v^3$
pada \cref{prop:b2:curves:kappaformula} adalah kesalahan yang paling sering
dalam praktik.
\end{remark}

\begin{figure}[ht]
\centering
\begin{tikzpicture}[scale=1.6]
  % parabola y = x^2/2
  \draw[->] (-2.2,0) -- (2.4,0) node[below] {$x$};
  \draw[->] (0,-0.4) -- (0,2.4) node[left] {$y$};
  \draw[thick, ocDarkBlue, domain=-2.1:2.1, smooth, samples=60]
    plot (\x, {0.5*\x*\x});
  % osculating circle at vertex: kappa = 1 => R = 1, center (0,1)
  \draw[ocRed, dashed] (0,1) circle (1);
  \fill[ocRed] (0,1) circle (0.03) node[right, font=\small]
    {pusat kelengkungan};
  \fill (0,0) circle (0.03);
  % tangent and normal at vertex
  \draw[->, thick] (0,0) -- (0.7,0) node[below, font=\small] {$T$};
  \draw[->, thick] (0,0) -- (0,0.7) node[left, font=\small] {$N$};
  % a second point
  \fill (1.2,0.72) circle (0.03);
  \draw[->, thick] (1.2,0.72) -- ++(0.448,0.5376)
    node[right, font=\small] {$T$};
  \draw[->, thick] (1.2,0.72) -- ++(-0.5376,0.448)
    node[above, font=\small] {$N$};
\end{tikzpicture}
\caption{Parabola $y = x^2/2$, kerangka Frenet bergeraknya $(T,
N)$, dan \omterm{def:b2:curves:curvature}{lingkaran oskulasi} di puncaknya (yang berjari-jari $1$, sebab
$\kappa(0) = 1$). Kerangkanya berputar ketika titiknya bergerak; jadi \omterm{thm:b2:curves:frenet2d}{kelengkungannya} adalah
laju perputaran itu per satuan \omterm{def:b2:curves:length}{panjang busur}.}
\label{fig:b2:curves:osculating}
\end{figure}

\begin{theorem}[Kelengkungan menentukan kurvanya]\label{thm:b2:curves:fundamental}
Misalkan $\kappa \colon J \to \R$ \omterm{def:b2:metric:continuity}{kontinu}. Maka ada sebuah busur
$\mathcal{C}^2$ berkelajuan satuan pada bidangnya dengan \omterm{thm:b2:curves:frenet2d}{kelengkungan}
$\kappa$, dan ia tunggal hingga sebuah isometri langsung (yakni perputaran
yang disusul penggeseran).
\end{theorem}

\begin{proof}
\emph{Keberadaannya.} Tetapkan $s_0 \in J$ lalu tetapkan $\varphi(s) =
\int_{s_0}^s \kappa(u)\,\dd u$, kemudian
\[
\tilde\gamma(s) = \Bigl(\int_{s_0}^s \cos\varphi(u)\,\dd u,\
\int_{s_0}^s \sin\varphi(u)\,\dd u\Bigr).
\]
Maka $T(s) = (\cos\varphi(s), \sin\varphi(s))$ merupakan vektor satuan,
$N(s) = (-\sin\varphi, \cos\varphi)$, dan
\[
T'(s) = \varphi'(s)\,(-\sin\varphi, \cos\varphi)
= \kappa(s)\,N(s) :
\]
jadi busurnya berkelajuan satuan dengan \omterm{thm:b2:curves:frenet2d}{kelengkungan} $\kappa$.

\emph{Ketunggalannya.} Misalkan $\gamma_1, \gamma_2$ busur berkelajuan satuan dengan
\omterm{thm:b2:curves:frenet2d}{kelengkungan} yang sama. Tiap singgung satuannya terangkat menjadi sebuah fungsi
sudut, lewat konstruksi yang gamblang: pandanglah $T_j$ sebagai
bilangan kompleks $z_j = a_j + \iu b_j$ bermodulus $1$, pilihlah
$\varphi_j(0)$ dengan $z_j(0) = \eu^{\iu\varphi_j(0)}$, lalu tetapkan
\[
\varphi_j(s) = \varphi_j(0) + \int_0^s\det\bigl(T_j,
T_j'\bigr)(u)\,\dd u .
\]
Dari $\abs{z_j} = 1$: berlaku $\operatorname{Re}(\conj{z_j}\,z_j') =
0$, jadi $\conj{z_j}\,z_j' = \iu\det(T_j, T_j') =
\iu\,\varphi_j'$, yakni $z_j' = \iu\varphi_j'z_j$; lalu
\[
\bigl(z_j\,\eu^{-\iu\varphi_j}\bigr)'
= \eu^{-\iu\varphi_j}\bigl(z_j' - \iu\varphi_j'z_j\bigr) = 0,
\]
jadi $z_j = \eu^{\iu\varphi_j}$ di seluruhnya: sehingga $T_j =
(\cos\varphi_j, \sin\varphi_j)$ dengan $\varphi_j$ berkelas
$\mathcal C^1$. Lebih jauh $\det(T_j, T_j') = \det(T_j,
\kappa N_j) = \kappa$, jadi $\varphi_j' = \kappa$. Karena itu $\varphi_2 = \varphi_1 + c$ untuk sebuah
konstanta $c$: jadi $T_2$ adalah $T_1$ yang diputar sebesar sudut tetap $c$, sehingga
setelah diintegralkan, $\gamma_2 = \rho(\gamma_1) + w$ dengan $\rho$ menyatakan
perputaran bersudut $c$ dan $w$ sebuah vektor tetap.
\end{proof}

\begin{remark}
Inilah prototipe berdimensi satu bagi sebuah \emph{teorema
dasar geometri}: bahwa himpunan \omterm{def:b2:metric:complete}{lengkap} invarian lokalnya (di sini, satu
fungsi) menggolongkan objeknya hingga gerak kaku. Adapun versi
berdimensi tiganya di bawah menuntut dua invarian.
\end{remark}

\begin{example}[Kelengkungan tetap berarti
lingkaran]\label{ex:b2:curves:constantkappa}
Ambillah $\kappa \equiv \kappa_0 > 0$ pada rumus keberadaannya:
maka $\varphi(s) = \kappa_0 s$ dan
\[
\tilde\gamma(s) = \Bigl(\frac{\sin\kappa_0s}{\kappa_0},\
\frac{1 - \cos\kappa_0s}{\kappa_0}\Bigr) :
\]
yakni lingkaran berjari-jari $1/\kappa_0$ yang berpusat di $(0,
1/\kappa_0)$, dan ditempuh pada kelajuan satuan. Berkat paruh ketunggalan
teoremanya, \emph{setiap} busur berkelajuan satuan yang \omterm{thm:b2:curves:frenet2d}{kelengkungannya} tetap
$\kappa_0$ merupakan keping sebuah lingkaran berjari-jari
$1/\kappa_0$ (atau sebuah garis bila $\kappa_0 = 0$) --- yakni konvers
perhitungan pada \cref{ex:b2:curves:kappaexamples}, sekaligus
kasus bidang bagi \cref{exo:b2:curves:9}.
\end{example}

\begin{example}[Membangun ulang sebuah kurva dari
kelengkungannya]\label{ex:b2:curves:reconstruction}
Kurva berkelajuan satuan manakah yang berjari-jari \omterm{thm:b2:curves:frenet2d}{kelengkungan} $R(s) = 1 +
s^2$? Dengan mengikuti bukti keberadaannya dengan $\kappa(s) =
\frac1{1+s^2}$ dan $s_0 = 0$: berlaku $\varphi(s) = \arctan s$, jadi
\[
T(s) = (\cos\arctan s,\ \sin\arctan s)
= \Bigl(\frac{1}{\sqrt{1+s^2}},\
\frac{s}{\sqrt{1+s^2}}\Bigr),
\]
lalu setelah diintegralkan,
\[
\tilde\gamma(s) = \Bigl(\ln\bigl(s + \sqrt{1 + s^2}\bigr),\
\sqrt{1 + s^2} - 1\Bigr).
\]
Dengan menetapkan $x = \ln(s + \sqrt{1+s^2})$, yakni $s = \sinh x$,
koordinat keduanya menjadi $\cosh x - 1$: jadi kurvanya adalah
\emph{katenari} $y = \cosh x - 1$. Ini menutup lingkarannya dengan
\cref{exo:b2:curves:3}, tempat kita menghitung $R = \cosh^2 x = 1 +
\sinh^2 x = 1 + s^2$ secara langsung: sebab teorema dasarnya
menjamin bahwa katenari adalah \emph{satu-satunya} kurva dengan profil
\omterm{thm:b2:curves:frenet2d}{kelengkungan} ini, hingga sebuah isometri langsung.
\end{example}

\begin{remark}[Di mana kelengkungan dipakai berikutnya]
Penguraian $\gamma'' = v'T + v^2\kappa N$ yang diperoleh pada
bukti \cref{prop:b2:curves:kappaformula} adalah
kinematika setiap gerak yang melengkung: yakni percepatan tangensial lawan
percepatan sentripetal. Adapun \omterm{thm:b2:curves:frenet2d}{kelengkungannya} kembali bagi permukaan
(\cref{ch:b2:surfaces}) lewat \omterm{thm:b2:curves:frenet2d}{kelengkungan} kurva yang digambar
padanya, dan kalkulus \omterm{pb:b2:curves:1}{selubung} pada soal akhir pekan bab ini
--- yakni evolut dan kaustik --- adalah optika geometris bagi
muka gelombang. Sedangkan jilid Tahun ke-3 mengangkat kembali sudut pandang yang
hakiki itu bagi submanifold $\R^n$.
\end{remark}

\section{Kerangka Frenet di ruang}

Kini misalkan $\tilde\gamma \colon J \to \R^3$ busur
$\mathcal{C}^3$ berkelajuan satuan yang \emph{biregular}: yakni $T'(s) \neq 0$ untuk
setiap $s$. Maka $\kappa(s) = \norm{T'(s)} > 0$ mendefinisikan
\emph{\omterm{thm:b2:curves:frenet2d}{kelengkungannya}} (yang tak bertanda di ruang: sebab tak ada orientasi normal
yang diistimewakan), lalu kita tetapkan:
\[
N(s) = \frac{T'(s)}{\kappa(s)}
\quad\text{(normal utama)},
\qquad
B(s) = T(s) \wedge N(s)
\quad\text{(binormal)} ,
\]
sehingga $(T, N, B)$ merupakan kerangka ortonormal yang langsung, yakni
\emph{kerangka Frenet}. Adapun bidang lewat $\tilde\gamma(s)$ yang direntang
$T, N$ disebut \emph{bidang oskulasi}.

\begin{theorem}[Rumus Frenet di ruang]\label{thm:b2:curves:frenet3d}
Ada fungsi \omterm{def:b2:metric:continuity}{kontinu} $\tau$, yakni \emph{torsi}\index{torsi}, dengan
\[
T' = \kappa N, \qquad
N' = -\kappa T + \tau B, \qquad
B' = -\tau N .
\]
\end{theorem}

\begin{proof}
Rumus pertamanya adalah definisi $N$. Adapun tiap vektor
$T, N, B$ bernorma tetap $1$ dan ketiganya ortogonal berpasangan;
jadi menurunkan keenam hubungan $\langle X, Y\rangle = \delta_{XY}$
menunjukkan bahwa matriks $(T', N', B')$ dalam basis $(T, N, B)$ bersifat
antisimetrik: sebab memang $\langle X', Y\rangle + \langle X, Y'\rangle = 0$ dan
$\langle X', X\rangle = 0$. Adapun entri $(N, T)$-nya adalah $\langle N', T\rangle =
-\langle N, T'\rangle = -\kappa$, sedangkan entri kolom $(T, B)$-nya
$\langle T', B\rangle = \kappa\langle N, B\rangle = 0$. Lalu menamai entri
bebas sisanya $\tau = \langle N', B\rangle$ memberi persis ketiga rumus pada
displainya: sebab keantisimetrikannya mengisi $\langle B', N\rangle = -\tau$ dan
$\langle B', T\rangle = 0$. Adapun kekontinuan $\tau = \langle N', B\rangle$ sudah jelas
karena $N'$ dan $B$ \omterm{def:b2:metric:continuity}{kontinu} (sebab $\tilde\gamma$ bersifat
$\mathcal{C}^3$, jadi $N = T'/\kappa$ bersifat $\mathcal{C}^1$).
\end{proof}

\begin{example}[Vektor Darboux]\label{ex:b2:curves:darboux}
Ketiga rumus Frenetnya memampat menjadi satu. Tetapkan $\omega(s) =
\tau\,T + \kappa\,B$ (yakni \emph{vektor Darboux}). Dengan memakai $B
\wedge T = N$, $T \wedge N = B$, dan $N \wedge B = T$:
\[
\omega \wedge T = \kappa\,N = T', \qquad
\omega \wedge N = \tau\,B - \kappa\,T = N', \qquad
\omega \wedge B = -\tau\,N = B' :
\]
jadi tiap vektor kerangkanya berevolusi lewat $X' = \omega \wedge X$, yakni
tanda kinematik sebuah \emph{perputaran seketika} dengan
vektor kecepatan sudut $\omega$. Kerangkanya berpusing berlaju
$\norm\omega = \sqrt{\kappa^2 + \tau^2}$ terhadap sumbu bergerak
$\omega$; jadi \omterm{thm:b2:curves:frenet2d}{kelengkungannya} adalah komponen pusingannya terhadap
binormalnya, sedangkan \omterm{thm:b2:curves:frenet3d}{torsinya} komponen terhadap singgungnya. Untuk
heliksnya, $\omega$ merupakan vektor tetap sepanjang sumbu
silindernya --- dan persis itulah sebabnya kerangka heliksnya berpresesi
merata. Adapun keantisimetrikan matriks Frenetnya, yang dieksploitasi pada
\cref{exo:b2:curves:7}, adalah bentuk matriks bagi satu
fakta geometris ini.
\end{example}

\begin{proposition}[Torsi mengukur kebidangan]\label{prop:b2:curves:torsion}
Sebuah busur biregular termuat pada sebuah bidang bila dan hanya bila $\tau
\equiv 0$; dan pada kasus itu bidangnya adalah bidang oskulasinya (yang
tetap).
\end{proposition}

\begin{proof}
Bila $\tau \equiv 0$, maka $B' = 0$, jadi $B$ merupakan vektor satuan tetap
$B_0$, dan
\[
\frac{\dd}{\dd s}\langle \tilde\gamma(s), B_0\rangle
= \langle T(s), B_0\rangle = 0 :
\]
jadi $\langle \tilde\gamma, B_0\rangle$ bersifat tetap, sehingga busurnya terletak pada bidang
yang ortogonal terhadap $B_0$. Sebaliknya, bila busurnya terletak pada bidang $P$,
maka $T$ dan $T'$ (sehingga $N$) sejajar dengan arah
$P$ untuk setiap $s$; jadi $B = T \wedge N$ adalah salah satu dari kedua normal
satuan $P$, dan karena ia \omterm{def:b2:metric:continuity}{kontinu} maka ia tetap; lalu $0 = B'
= -\tau N$ dengan $N \neq 0$ memaksa $\tau \equiv 0$.
\end{proof}

\begin{example}[Sebuah lingkaran yang miring bertorsi nol]
Busur $\gamma(t) = \bigl(\cos t,\ \tfrac{\sin t}{\sqrt2},\
\tfrac{\sin t}{\sqrt2}\bigr)$ terletak pada bidang $y = z$, dan
merupakan lingkaran satuan bidang itu (periksa: sebab $\norm{\gamma(t)} =
1$ dan basis ortonormal bidangnya $(1,0,0)$ dan
$(0, \tfrac1{\sqrt2}, \tfrac1{\sqrt2})$ memamerkan parameterisasi
bakunya). Tanpa satu perhitungan Frenet pun,
\cref{prop:b2:curves:torsion} meramalkan $\tau \equiv 0$, dan
binormal tetapnya mestilah normal satuan bidangnya $\pm(0,
\tfrac1{\sqrt2}, -\tfrac1{\sqrt2})$. Jadi \omterm{thm:b2:curves:frenet3d}{torsi} tak mengukur
keadaan ``miring di ruang''; ia mengukur \emph{meninggalkan} sebuah
bidang. Hanya $b$ yang tak nol pada heliks di bawah yang menghasilkan \omterm{thm:b2:curves:frenet3d}{torsi}
yang sejati.
\end{example}

\begin{example}[Heliksnya]\label{ex:b2:curves:helixfrenet}
Untuk heliks $\gamma(t) = (a\cos t, a\sin t, bt)$ dengan $a > 0$, kita telah
menghitung $s = ct$ dengan $c = \sqrt{a^2 + b^2}$. Maka
\[
T = \frac1c(-a\sin t,\ a\cos t,\ b),
\qquad
T' \cdot \frac{\dd t}{\dd s}
= \frac{1}{c^2}(-a\cos t, -a\sin t, 0),
\]
jadi $\kappa = a/c^2 = a/(a^2 + b^2)$ dan $N = (-\cos t, -\sin t,
0)$: sehingga normal utamanya menunjuk mendatar ke arah sumbunya.
Berikutnya $B = T \wedge N = \frac1c(b\sin t, -b\cos t, a)$, dan $B'
\frac{\dd t}{\dd s} = \frac{b}{c^2}(\cos t, \sin t, 0) = -\tau N$
memberi
\[
\boxed{\ \kappa = \frac{a}{a^2 + b^2}, \qquad
\tau = \frac{b}{a^2 + b^2}. \ }
\]
Kedua invariannya bersifat tetap --- dan sebaliknya dapat ditunjukkan bahwa
satu-satunya kurva biregular dengan $\kappa > 0$ tetap dan $\tau$ yang
tetap adalah heliks (yakni lingkaran ketika $\tau = 0$). Perhatikan tandanya: $b >
0$ memberi heliks tangan-kanan dengan \omterm{thm:b2:curves:frenet3d}{torsi} yang positif.
\end{example}

\begin{example}[Kelengkungan dan torsi tanpa panjang busur; yakni
kubik terpuntir]\label{ex:b2:curves:twistedcubic}
Memparameterkan ulang dengan \omterm{def:b2:curves:length}{panjang busur} biasanya mustahil dalam bentuk
tertutup, jadi invariannya harus disarikan dari turunan
mentahnya. Tulislah $v = \norm{\gamma'} = s'$; maka $\gamma' =
vT$ dan, seperti pada bukti
\cref{prop:b2:curves:kappaformula},
\[
\gamma'' = v'T + v^2\kappa N,
\qquad
\gamma' \wedge \gamma'' = v^3\kappa\,(T \wedge N) =
v^3\kappa\,B .
\]
Setelah \omterm{def:b2:nvs:norm}{normanya} diambil (dengan $\kappa \geq 0$ di ruang):
\[
\kappa = \frac{\norm{\gamma' \wedge \gamma''}}{v^3} .
\]
Lalu menurunkan $\gamma''$ sekali lagi dan mengubah $N' =
v(-\kappa T + \tau B)$ (lewat aturan rantai melalui $s$), satu-satunya
komponen-$B$-nya datang dari suku terakhirnya:
\[
\gamma''' = \bigl(v'' - v^3\kappa^2\bigr)T +
\bigl(v'v\kappa + (v^2\kappa)'\bigr)N + v^3\kappa\tau\,B,
\]
sehingga, setelah dipasangkan dengan $\gamma' \wedge \gamma'' = v^3\kappa B$,
\[
\det(\gamma', \gamma'', \gamma''') = \langle\gamma' \wedge
\gamma'',\ \gamma'''\rangle = v^6\kappa^2\tau,
\qquad\text{yakni}\qquad
\tau = \frac{\det(\gamma', \gamma'',
\gamma''')}{\norm{\gamma' \wedge \gamma''}^2} .
\]
Penerapannya pada \emph{kubik terpuntir} $\gamma(t) = (t,\
t^2,\ t^3)$ di $t = 0$: berlaku $\gamma' = (1, 0, 0)$, $\gamma'' =
(0, 2, 0)$, $\gamma''' = (0, 0, 6)$, jadi $v = 1$,
\[
\gamma' \wedge \gamma'' = (0, 0, 2), \qquad \kappa(0) = 2,
\qquad
\det(\gamma', \gamma'', \gamma''') = 12, \qquad \tau(0) =
\frac{12}{4} = 3 .
\]
Pelajaran penutupnya: kedua rumusnya adalah nisbah yang di dalamnya kelajuan
$v$ tercoret persis sejauh yang diperlukan --- sebab $\kappa$ berskala
seperti turunan kedua untuk tiap satuan \omterm{def:b2:curves:length}{panjang}, sedangkan $\tau$ seperti
volume campuran tiga turunan per luas kuadrat --- dan itulah
\emph{sebabnya} keduanya geometris sedangkan $\gamma''$ sendiri
tidak.
\end{example}

\begin{remark}[Teorema dasar bagi kurva ruang]
Seperti pada bidangnya, pasangan $(\kappa, \tau)$ dengan $\kappa > 0$
menentukan sebuah busur biregular hingga isometri langsung $\R^3$: sebab
rumus Frenetnya membentuk sistem \omterm{def:b2:diffcalc:differential}{diferensial} linear bagi kerangkanya
$(T, N, B)$, yang kepadanya teori Cauchy--Lipschitz pada
\cref{ch:b2:diffeq} berlaku; lalu keortonormalan kerangka penyelesaiannya
terpelihara karena matriks koefisiennya bersifat antisimetrik (yakni hujah
matriks Gram yang sama seperti pada \cref{exo:b2:curves:7}), dan kurvanya
dipulihkan dengan mengintegralkan $T$. Kami serahkan rinciannya kepada
pembaca sebagai latihan yang berat tetapi mendidik.
\end{remark}

\section{Kajian lokal: kedudukan terhadap garis singgungnya}

\begin{proposition}[Bentuk lokal di sebuah titik
regular]\label{prop:b2:curves:local}
Misalkan $\gamma$ busur bidang berkelas $\mathcal{C}^k$ di $t_0$, dengan
$p$ indeks terkecil yang $\gamma^{(p)}(t_0) \neq 0$ dan $q$
indeks terkecil $> p$ yang $\gamma^{(q)}(t_0)$ tak segaris dengan
$\gamma^{(p)}(t_0)$ (dengan mengandaikan keduanya ada, $q \leq k$). Dalam basis
$(u, v) = (\gamma^{(p)}(t_0), \gamma^{(q)}(t_0))$ yang dipusatkan di
$\gamma(t_0)$, Taylor--Young memberi koordinat
\[
X(t) \sim \frac{(t - t_0)^p}{p!},
\qquad
Y(t) \sim \frac{(t - t_0)^q}{q!} .
\]
Adapun gambaran lokalnya hanya bergantung pada keparitasan $p$ dan $q$:
\begin{center}
\begin{tabular}{c c l}
$p$ ganjil, $q$ genap & \emph{titik biasa} & tetap di satu sisi
singgungnya\\
$p$ ganjil, $q$ ganjil & \emph{titik belok} & kurvanya menyeberangi
singgungnya\\
$p$ genap, $q$ ganjil & \emph{titik runcing jenis pertama} & kedua
cabangnya di sisi yang berlawanan\\
$p$ genap, $q$ genap & \emph{titik runcing jenis kedua} & kedua
cabangnya di sisi yang sama\\
\end{tabular}
\end{center}
\end{proposition}

\begin{proof}
Taylor--Young pada orde $q$ (sebab fungsi $\gamma$ bersifat
$\mathcal{C}^q$ di dekat $t_0$):
\[
\gamma(t) - \gamma(t_0)
= \sum_{j=p}^{q} \frac{(t-t_0)^j}{j!}\,\gamma^{(j)}(t_0)
+ o\bigl((t-t_0)^q\bigr).
\]
Berkat pemilihan $p$ dan $q$, tiap $\gamma^{(j)}(t_0)$ dengan $p
\leq j < q$ segaris dengan $u$; jadi setelah mengumpulkan komponennya dalam
basis $(u, v)$: $X(t) = \frac{(t-t_0)^p}{p!}(1 + o(1))$ dan
$Y(t) = \frac{(t-t_0)^q}{q!}(1 + o(1))$. Adapun tabel tanda $X$ dan
$Y$ untuk $t \gtrless t_0$ --- yang persis dikendalikan keparitasannya ---
memberi keempat gambarannya: misalnya bila $p$ genap, maka $X > 0$ pada
kedua sisinya (sebab kedua cabangnya pergi ke arah $+u$: yakni titik runcing),
dan sisi \omterm{def:b2:curves:arc}{garis singgungnya} ($\operatorname{sign} Y$) berbalik
ketika $q$ ganjil.
\end{proof}

\begin{example}
Untuk $\gamma(t) = (t^2, t^3)$ di $t_0 = 0$
(\cref{ex:b2:curves:cusp}): berlaku $\gamma'' (0)= (2, 0)$ dan $\gamma'''(0) =
(0, 6)$, jadi $p = 2$ dan $q = 3$: yakni titik runcing jenis pertama, yaitu
gambaran parabola semikubik yang sudah dikenal. Sedangkan untuk $\gamma(t) = (t,
t^3)$ di $0$: $p = 1$ dan $q = 3$: yakni titik belok --- sebab kubiknya menyeberangi
singgungnya.
\end{example}

\begin{remark}[Pandangan ke depan di dalam jilid ini]
Kurva menyuapi bab berikutnya lewat tiga cara. Bila digambar pada sebuah
permukaan, ia mendefinisikan bidang singgung dan bentuk fundamental
pertamanya (\cref{ch:b2:surfaces}), dan \omterm{def:b2:curves:length}{panjangnya}
dihitung dengan membatasi metrik ruang lingkupnya --- jadi
bab berikutnya sebagian besar adalah bab ini yang direlatifkan. Adapun
kalkulus \omterm{pb:b2:curves:1}{selubung} pada soal akhir pekannya bertemu integral
lipat pada \cref{ch:b2:multint}, tempat luas \omterm{pb:b2:curves:1}{astroidnya}
dihitung ulang lewat rumus Green
(\cref{exo:b2:multint:5}) --- yakni satu kurva, dua teori, dan
jawaban yang cocok. Lalu sistem Frenetnya sudah memakai
persamaan \omterm{def:b2:diffcalc:differential}{diferensial} linear pada \cref{ch:b2:diffeq}
(yakni keberadaannya, ketunggalannya, dan hujah pemeliharaan keortogonalan
pada \cref{exo:b2:curves:7}): jadi teorema dasar
kurva adalah teorema persamaan \omterm{def:b2:diffcalc:differential}{diferensial} yang mengenakan
busana geometri.
\end{remark}

\begin{remark}[Metode: menjalankan kajian lokalnya]
Dalam praktik penggolongannya berupa rutin empat langkah.
\emph{Satu}, turunkan di $t_0$ sampai turunan tak nol
pertamanya muncul: indeksnya adalah $p$, nilainya vektor
$u$. \emph{Dua}, teruskanlah menurunkan sampai sebuah turunan
yang tak segaris dengan $u$ muncul: indeksnya $q$, vektornya $v$.
\emph{Tiga}, bacalah keparitasan $(p, q)$ pada tabelnya.
\emph{Empat}, gambarlah: kurvanya pergi sepanjang $+u$ bila $p$ ganjil
(atau sepanjang $u$ lalu kembali sepanjang $u$ bila $p$ genap), pada sisi
$v$ yang didiktekan tanda $Y$. Ada dua peringatan. Kerangka
$(u, v)$ umumnya \emph{tak} ortonormal --- sebab tabelnya
memaparkan kedudukan relatif terhadap \omterm{def:b2:curves:arc}{garis singgungnya}, bukan sudut
atau jarak, jadi janganlah membaca \omterm{thm:b2:curves:frenet2d}{kelengkungan} dari gambarnya.
Lalu turunan antara yang segaris dengan $u$ memang diizinkan
di antara pangkat $p$ dan $q$ (sebab ia hanya menggeser uraian
$X$); yang \emph{tak} boleh terjadi adalah berhenti pada turunan
tak nol pertamanya lalu menebak $q = p + 1$: sebab untuk $\gamma(t)
= (t^2, t^4 + t^5)$ tebakan naif $q = 3$ keliru, karena
$\gamma^{(3)}(0)$ masih segaris dengan $\gamma''(0)$ ---
dan inilah persis \cref{exo:b2:curves:5}.
\end{remark}

\section{Latihan}

\begin{exercise}[$\star$]\label{exo:b2:curves:1}
Hitunglah \omterm{def:b2:curves:length}{panjang} satu lengkung sikloid $\gamma(t) = (t -
\sin t,\ 1 - \cos t)$ dengan $t \in [0, 2\pi]$. \emph{(Pakailah $1 - \cos t =
2\sin^2(t/2)$.)}
\end{exercise}

\begin{exercise}[$\star$]\label{exo:b2:curves:2}
Hitunglah \omterm{thm:b2:curves:frenet2d}{kelengkungan} elips $\gamma(t) = (a\cos t,\ b\sin
t)$ (dengan $a > b > 0$) lalu tentukan letak titik yang \omterm{thm:b2:curves:frenet2d}{kelengkungannya}
maksimum dan minimum.
\end{exercise}

\begin{exercise}[$\star$]\label{exo:b2:curves:3}
Tunjukkan bahwa \omterm{def:b2:curves:length}{panjang busur} grafik $f(x) = \cosh x$ pada
$[0, x]$ sama dengan $\sinh x$, lalu hitunglah \omterm{thm:b2:curves:frenet2d}{kelengkungan} kurva
ini (yakni \emph{katenari}). Periksalah bahwa $R(x) = 1/\kappa(x) =
\cosh^2 x$.
\end{exercise}

\begin{exercise}[$\star\star$]\label{exo:b2:curves:4}
(Spiral logaritmik) Misalkan $\gamma(t) = e^{t}(\cos t,\ \sin t)$ dengan $t
\in \R$.
Tunjukkan bahwa sudut antara $\gamma(t)$ dan $\gamma'(t)$ bersifat
tetap, lalu hitunglah \omterm{def:b2:curves:length}{panjang busur} $\gamma$ pada $(-\infty, 0]$
(yang berhingga!), beserta \omterm{thm:b2:curves:frenet2d}{kelengkungannya}.
\end{exercise}

\begin{exercise}[$\star\star$]\label{exo:b2:curves:5}
Tentukan $p$, $q$ dan bentuk lokalnya (biasa, belok,
runcing) bagi $\gamma(t) = (t^2,\ t^4 + t^5)$ di $t = 0$, dan bagi
$\gamma(t) = (t^3,\ t^4)$ di $t = 0$.
\end{exercise}

\begin{exercise}[$\star\star$]\label{exo:b2:curves:6}
Misalkan $\gamma$ busur bidang berkelajuan satuan dengan $\kappa(s) > 0$ untuk
setiap $s$, lalu misalkan $c(s) = \gamma(s) + \frac{1}{\kappa(s)}N(s)$
\omterm{def:b2:curves:curvature}{pusat kelengkungannya} (yakni kurva $c$ disebut \emph{evolut}).
Dengan mengandaikan $\kappa$ bersifat $\mathcal{C}^1$, tunjukkanlah bahwa $c'(s) =
-\frac{\kappa'(s)}{\kappa(s)^2}N(s)$: jadi evolutnya bersinggungan dengan
garis normal $\gamma$.
\end{exercise}

\begin{exercise}[$\star\star\star$]\label{exo:b2:curves:7}
Misalkan $A(s)$ keluarga \omterm{def:b2:metric:continuity}{kontinu} matriks antisimetrik $3 \times 3$
dan $F' = F A$ sebuah penyelesaian matriks dengan $F(s_0)$
yang ortogonal. Tunjukkan bahwa $F(s)$ ortogonal untuk setiap $s$.
\emph{(Turunkan $G = F F^{\mathsf T}$ lalu pakailah ketunggalan pada
Cauchy--Lipschitz.)} Jelaskanlah kaitannya dengan sistem Frenet.
\end{exercise}

\begin{exercise}[$\star\star\star$]\label{exo:b2:curves:8}
(\omterm{thm:b2:curves:frenet2d}{Kelengkungan} total sebuah kurva tertutup yang \omterm{def:b2:affine:convex}{cembung}) Misalkan $\tilde\gamma$
busur bidang tertutup $\mathcal{C}^2$ berkelajuan satuan yang berpanjang $L$
(jadi $\tilde\gamma(s + L) = \tilde\gamma(s)$), yang ditempuh sekali
berlawanan arah jarum jam. Dengan memakai fungsi sudut $\varphi$ yang $T =
(\cos\varphi, \sin\varphi)$ dari
\cref{thm:b2:curves:fundamental}, jelaskanlah mengapa $\varphi(L) -
\varphi(0)$ merupakan kelipatan $2\pi$, lalu tunjukkan bahwa
$\int_0^L \kappa(s)\,\dd s = \varphi(L) - \varphi(0)$. (Untuk sebuah
lingkaran berjari-jari $R$: $\int \kappa = \frac1R \cdot 2\pi R = 2\pi$.
Adapun teorema singgung berputar menegaskan nilainya $2\pi$ bagi
setiap kurva tertutup sederhana; tetapi kamu tak diminta membuktikannya.)
\end{exercise}

\begin{exercise}[$\star\star\star$]\label{exo:b2:curves:9}
Tunjukkan bahwa kurva ruang biregular dengan $\kappa > 0$ yang tetap dan
$\tau = 0$ merupakan (sebuah busur) lingkaran berjari-jari $1/\kappa$.
\emph{(Pakailah \cref{prop:b2:curves:torsion}, lalu tunjukkan bahwa pusatnya
$\gamma + \frac1\kappa N$ bersifat tetap.)}
\end{exercise}

\begin{exercise}[$\star$]\label{exo:b2:curves:10}
Hitunglah \omterm{def:b2:curves:length}{panjang busur} parabola $y = x^2/2$ pada
$\intcc0a$ lalu tunjukkan bahwa ia sama dengan
\[
\tfrac12\Bigl(a\sqrt{1 + a^2} + \ln\bigl(a + \sqrt{1 +
a^2}\bigr)\Bigr).
\]
\end{exercise}

\begin{exercise}[$\star\star$]\label{exo:b2:curves:11}
Misalkan $\gamma$ busur $\mathcal C^2$ yang regular di $\R^n$ dan semua
\omterm{def:b2:curves:arc}{garis singgungnya} melalui titik tetap $P$. Buktikan bahwa
lintasan $\gamma$ termuat pada sebuah garis lurus.
\emph{(Parameterkan lewat \omterm{def:b2:curves:length}{panjang busur}, tulislah $\gamma(s) +
\lambda(s)T(s) = P$ lalu turunkan.)}
\end{exercise}

\begin{exercise}[$\star\star\star$]\label{exo:b2:curves:12}
(Teorema dasar bagi kurva ruang) Misalkan $\kappa > 0$ dan
$\tau$ fungsi \omterm{def:b2:metric:continuity}{kontinu} pada sebuah interval $J$. Kerjakanlah
program pada catatan yang menyusul
\cref{ex:b2:curves:helixfrenet}: (a) tunjukkan bahwa sistem
linear $F' = FA(s)$, dengan $A(s)$ matriks Frenet yang
antisimetrik yang \omterm{def:b2:structures:generated}{dibangun} dari $\kappa, \tau$ dan $F(s_0)$ sebuah
kerangka ortonormal langsung, punya satu penyelesaian global tunggal yang tetap
berupa kerangka ortonormal langsung; (b) bangunlah sebuah kurva biregular
berkelajuan satuan dengan \omterm{thm:b2:curves:frenet2d}{kelengkungan} $\kappa$ dan \omterm{thm:b2:curves:frenet3d}{torsi} $\tau$;
(c) buktikan ketunggalannya hingga sebuah isometri langsung $\R^3$.
\end{exercise}

\section{Soal: selubung --- astroid, dua evolut, dan
sebuah kaustik}

\begin{omfigure}
\begin{tikzpicture}[scale=2.2]
  \draw[->] (-1.15,0) -- (1.25,0) node[below] {$x$};
  \draw[->] (0,-1.15) -- (0,1.25) node[left] {$y$};
  \draw[omDef] (1,0) -- (0,0.0) (0.966,0) -- (0,0.259)
    (0.866,0) -- (0,0.5) (0.707,0) -- (0,0.707)
    (0.5,0) -- (0,0.866) (0.259,0) -- (0,0.966);
  \draw[omThm, very thick]
    plot[domain=0:360, samples=120]
    ({cos(\x)*cos(\x)*cos(\x)}, {sin(\x)*sin(\x)*sin(\x)});
\end{tikzpicture}

{\small Sebuah tangga berpanjang $1$ yang meluncur turun di dinding (yakni
kedudukan birunya) tak pernah menyeberangi \omterm{pb:b2:curves:1}{astroid} $x^{2/3} + y^{2/3} = 1$
(yang merah): jadi \omterm{pb:b2:curves:1}{astroidnya} adalah \emph{\omterm{pb:b2:curves:1}{selubung}} keluarga
ruas itu, yang bersinggungan dengan setiap ruasnya.}
\end{omfigure}

\begin{problem}[{Soal akhir pekan --- mesin selubung dan
empat kurva klasik}]\label{pb:b2:curves:1}
Sebuah keluarga garis berparameter satu biasanya gagal meliputi
bidangnya secara merata: sebab garisnya menumpuk sepanjang sebuah kurva yang bersinggungan dengan
semuanya, yakni \emph{selubung}\index{selubung} keluarga itu. Adapun sinar cahaya
membuat \omterm{pb:b2:curves:1}{selubungnya} kasatmata sebagai \emph{kaustik} --- yakni kurva
terang beruncing di dalam cangkir kopi. Soal ini membangun
mesin \omterm{pb:b2:curves:1}{selubung} yang umum, lalu menjalankannya empat kali: yakni pada tangga
yang meluncur (\omterm{pb:b2:curves:1}{astroid}), pada normal parabola dan
sikloid (evolut, dengan bandul Huygens di akhirnya), dan pada
kaustik cangkir kopi (\omterm{pb:b2:curves:1}{nefroid}). Di sepanjang soal ini, $D_t$ menyatakan
garis berpersamaan $a(t)\,x + b(t)\,y = c(t)$, dengan $a, b, c$
berupa fungsi $\mathcal C^2$ dan $(a(t), b(t)) \neq (0,0)$,
serta $\Delta(t) = a(t)b'(t) - a'(t)b(t)$.

\medskip
\noindent\textbf{Bagian I --- Mesin \omterm{pb:b2:curves:1}{selubungnya}.}
\begin{enumerate}
  \item Andaikan $\Delta(t) \neq 0$. Tunjukkan bahwa \emph{sistem
        karakteristiknya}
        \[
        \begin{cases} a(t)\,x + b(t)\,y = c(t)\\
        a'(t)\,x + b'(t)\,y = c'(t)\end{cases}
        \]
        punya satu penyelesaian tunggal $E(t) = (x(t), y(t))$, yang diberikan
        $x = \dfrac{cb' - c'b}{\Delta}$ dan $y = \dfrac{ac' -
        a'c}{\Delta}$.
  \item Andaikan lebih jauh bahwa $E$ bersifat $\mathcal C^1$ di dekat $t$
        dengan $E'(t) \neq 0$. Dengan menurunkan persamaan pertama
        sistemnya, tunjukkanlah $a(t)\,x'(t) +
        b(t)\,y'(t) = 0$, lalu simpulkan bahwa kurva $E$
        melalui sebuah titik $D_t$ \emph{dengan arah}
        $D_t$: jadi keluarganya bersinggungan dengan $E$,
        yang disebut \emph{\omterm{pb:b2:curves:1}{selubungnya}}.
  \item Pemeriksaan kewarasannya: \omterm{def:b2:curves:arc}{garis singgung} parabola $y =
        x^2/2$ di titik $(t, t^2/2)$ adalah $tx - y =
        t^2/2$. Periksalah bahwa mesin \omterm{pb:b2:curves:1}{selubungnya} mengembalikan
        parabolanya sendiri.
  \item (\omterm{pb:b2:curves:1}{Selubung} normalnya) Misalkan $\gamma$ berkelajuan satuan
        dengan $\kappa(s) \neq 0$. Adapun garis normal di
        $\gamma(s)$ adalah $\{M : \langle M - \gamma(s), T(s)
        \rangle = 0\}$. Tunjukkan bahwa sistem karakteristiknya
        memaksa $\langle M - \gamma(s), N(s)\rangle =
        1/\kappa(s)$, sehingga titik karakteristiknya adalah
        \omterm{def:b2:curves:curvature}{pusat kelengkungannya}: jadi \emph{\omterm{pb:b2:curves:1}{selubung}
        normalnya adalah evolutnya}, yang memulihkan
        \cref{exo:b2:curves:6}. Periksalah $\Delta(s) =
        \kappa(s)$.
  \item Dua kemerosotan. Untuk berkas $D_\theta : x\cos
        \theta + y\sin\theta = 0$, tunjukkan bahwa
        titik karakteristiknya adalah titik asalnya untuk setiap $\theta$
        (jadi ``\omterm{pb:b2:curves:1}{selubungnya}'' runtuh menjadi satu titik, dan $E' = 0$:
        sehingga pertanyaan 2 tak berlaku). Sedangkan untuk keluarga garis sejajar
        (dengan $a, b$ yang tetap), tunjukkanlah $\Delta \equiv 0$ dan
        bahwa sistem karakteristiknya secara umum
        tak serasi: jadi tak ada \omterm{pb:b2:curves:1}{selubung}.
\end{enumerate}

\medskip
\noindent\textbf{Bagian II --- Tangga yang meluncur dan
\omterm{pb:b2:curves:1}{astroidnya}.} Sebuah ruas berpanjang $1$ meluncur dengan satu ujungnya $P_t =
(\cos t, 0)$ di lantai dan yang lain $Q_t = (0, \sin t)$ di
dinding, dengan $t \in \intoo0{\pi/2}$.
\begin{enumerate}[resume]
  \item Tunjukkan bahwa garis $(P_tQ_t)$ berpersamaan $x\sin t +
        y\cos t = \sin t\cos t$, dan bahwa mesin \omterm{pb:b2:curves:1}{selubungnya}
        memberi titik karakteristik
        \[
        E(t) = (\cos^3 t,\ \sin^3 t) :
        \]
        yakni \emph{astroid}\index{astroid}, yang berpersamaan
        implisit $x^{2/3} + y^{2/3} = 1$ (dan diperluas ke
        kuadran lainnya lewat kesimetrian).
  \item Tunjukkan bahwa $E'(0) = 0$ dan, dengan memakai penggolongan
        lokalnya (\cref{prop:b2:curves:local}), bahwa
        \omterm{pb:b2:curves:1}{astroidnya} punya titik runcing jenis pertama di $(1, 0)$
        --- dan demikian pula di keempat titik sumbunya.
  \item Hitunglah $\norm{E'(t)} = \tfrac32\abs{\sin 2t}$ lalu
        turunkan bahwa \omterm{def:b2:curves:length}{panjang} total \omterm{pb:b2:curves:1}{astroidnya} adalah $6$.
  \item Di mana tangganya menyentuh \omterm{pb:b2:curves:1}{astroidnya}? Tunjukkan $E(t) =
        P_t + \sin^2 t\,(Q_t - P_t)$: jadi titik sentuhnya
        membagi tangganya dengan nisbah $\sin^2 t : \cos^2
        t$, sambil menyapunya dari satu ujung ke ujung lain ketika
        tangganya meluncur.
  \item Hitunglah luas yang dilingkupi \omterm{pb:b2:curves:1}{astroidnya}: tunjukkan bahwa
        luas kuadran pertamanya adalah $3\int_0^{\pi/2}\sin^4
        t\cos^2 t\,\dd t$, nilailah integralnya lewat
        pelinearan (yakni $\sin^2 2t = \tfrac{1 - \cos 4t}2$),
        lalu simpulkan bahwa luas totalnya adalah $3\pi/8$.
\end{enumerate}

\medskip
\noindent\textbf{Bagian III --- Evolut parabolanya.}
Misalkan $\gamma(t) = (t, t^2/2)$.
\begin{enumerate}[resume]
  \item Tunjukkan bahwa garis normal di $\gamma(t)$ berpersamaan
        $x + t\,y = t + t^3/2$.
  \item Jalankan mesin \omterm{pb:b2:curves:1}{selubungnya}: tunjukkan bahwa \omterm{pb:b2:curves:1}{selubung}
        normalnya adalah
        \[
        E(t) = \Bigl(-t^3,\ 1 + \tfrac32 t^2\Bigr),
        \]
        dengan persamaan implisit $x^2 = \tfrac8{27}(y - 1)^3$:
        yakni sebuah parabola semikubik.
  \item Periksa silang dengan pertanyaan 4: hitunglah \omterm{def:b2:curves:curvature}{pusat
        kelengkungannya} $\gamma(t) + \frac1{\kappa(t)}N(t)$ dari
        $\kappa(t) = (1 + t^2)^{-3/2}$
        (\cref{ex:b2:curves:kappaexamples}) lalu perolehlah
        titik yang sama.
  \item Tunjukkan bahwa evolutnya punya titik runcing jenis pertama di
        $(0, 1)$, yakni \omterm{def:b2:curves:curvature}{pusat kelengkungan} di puncaknya ---
        yaitu titik tempat $\kappa$ ekstrem, sebagaimana diramalkan
        rumus $c' = -\frac{\kappa'}{\kappa^2}N$ pada
        \cref{exo:b2:curves:6}.
  \item Berapa banyak normal parabolanya yang melalui sebuah titik
        $(x_0, y_0)$? Tunjukkan bahwa jawabannya dikendalikan
        kubik $\tfrac{t^3}2 + (1 - y_0)\,t - x_0 = 0$;
        tanganilah kasus sumbunya $x_0 = 0$ selengkapnya (yakni satu normal
        bila $y_0 < 1$, dan tiga bila $y_0 > 1$), lalu tafsirkanlah
        evolutnya sebagai kurva peralihannya.
\end{enumerate}

\medskip
\noindent\textbf{Bagian IV --- Kaustik cangkir kopi.} Sinar
sejajar berarah $(1, 0)$ menghantam bagian dalam cermin
lingkaran $x^2 + y^2 = 1$; lalu sinar yang mengenai $P_\theta =
(\cos\theta, \sin\theta)$ terpantul menurut hukum
pemantulan.
\begin{enumerate}[resume]
  \item Dari kesimetrian cermin pada normalnya (yakni jari-jarinya),
        benarkanlah bahwa arah pantulnya adalah $v = u -
        2\langle u, n\rangle n$ dengan $u = (1,0)$ dan $n =
        (\cos\theta, \sin\theta)$, lalu hitunglah $v =
        -(\cos2\theta, \sin2\theta)$.
  \item Tunjukkan bahwa sinar pantulnya terletak pada garis
        \[
        x\sin 2\theta - y\cos 2\theta = \sin\theta .
        \]
  \item Jalankan mesin \omterm{pb:b2:curves:1}{selubungnya} (dengan $\Delta = 2$): tunjukkan bahwa
        kaustiknya adalah
        \[
        E(\theta) = \Bigl(\tfrac{3\cos\theta -
        \cos3\theta}4,\ \tfrac{3\sin\theta -
        \sin3\theta}4\Bigr),
        \]
        yakni \emph{nefroid}\index{nefroid}.
  \item Hitunglah $E'(\theta) = \tfrac32\sin\theta\,
        (\cos2\theta, \sin2\theta)$; periksalah bahwa arah
        singgungnya adalah arah sinar pantulnya (pertanyaan
        16), tentukan letak kedua titik runcingnya $(\pm\tfrac12, 0)$, lalu
        tunjukkan bahwa sinar pantulnya menyeberangi sumbu $y = 0$
        di $x = \frac1{2\cos\theta}$ --- jadi sinar yang hampir sesumbu
        memfokus di $x = \tfrac12$: yakni \omterm{def:b2:curves:length}{panjang} fokus $R/2$ bagi
        cermin berjari-jari $R$.
  \item Tunjukkan bahwa \omterm{pb:b2:curves:1}{nefroidnya} berpanjang total $6$ dan bahwa
        di dekat $\theta = 0$,
        \[
        E(\theta) - \bigl(\tfrac12, 0\bigr) =
        \bigl(\tfrac34\theta^2 + o(\theta^2),\ \theta^3 +
        o(\theta^3)\bigr) :
        \]
        yakni sebuah titik runcing jenis pertama, yang menunjuk sepanjang sumbunya.
  \item Jelaskanlah dalam satu paragraf mengapa kaustiknya terang:
        sebab lewat setiap titik tepat di luar kaustiknya melintas dua
        sinar pantul, sedangkan lewat setiap titik padanya sinarnya
        ``terpusat tanpa batas'' (karena pemetaan $(\theta,
        \text{jarak sepanjang sinar}) \mapsto \R^2$ punya
        titik kritis persis pada \omterm{pb:b2:curves:1}{selubungnya}).
\end{enumerate}

\medskip
\noindent\textbf{Bagian V --- Huygens: sikloid adalah evolutnya
sendiri.} Misalkan $\gamma(t) = (t - \sin t,\ 1 - \cos t)$ dengan $t \in
\intoo0{2\pi}$, yakni satu lengkung sikloid.
\begin{enumerate}[resume]
  \item Hitunglah $\kappa(t) = -\dfrac1{4\sin(t/2)}$ beserta
        \omterm{def:b2:curves:curvature}{pusat kelengkungannya}; lalu tunjukkan bahwa evolutnya adalah
        \[
        c(t) = (t + \sin t,\ \cos t - 1),
        \]
        dan bahwa penyulihan $t = u + \pi$ memamerkannya
        sebagai sikloid semula yang digeser sebesar $(\pi, -2)$:
        jadi \emph{evolut sebuah sikloid adalah sikloid yang
        kongruen} (yakni Huygens).
  \item Periksalah bahwa \omterm{def:b2:curves:curvature}{jari-jari kelengkungan} di puncaknya $t =
        \pi$ sama dengan $4$, yakni separuh \omterm{def:b2:curves:length}{panjang} $8$ satu
        lengkungnya (\cref{exo:b2:curves:1}); lalu tentukan letak titik runcing
        evolutnya tepat di bawah puncaknya, pada jarak $4$.
  \item (Sifat talinya) Misalkan $\gamma$ berkelajuan satuan dengan
        $\kappa > 0$, dengan $\kappa$ berkelas $\mathcal C^1$ dan
        $R = 1/\kappa$ monoton tegas. Dengan memakai $c' = R'N$,
        tunjukkanlah bahwa \omterm{def:b2:curves:length}{panjang busur} evolutnya antara
        $c(s_0)$ dan $c(s_1)$ adalah $\abs{R(s_1) - R(s_0)}$.
        Tafsirkanlah: sebuah tali tegang yang terurai dari evolutnya,
        yang berpanjang $R(s_0)$ pada awalnya, membuat ujung bebasnya
        melacak kurva semula --- jadi sebuah bandul yang berayun
        di antara dua pipi sikloid pada jam Huygens
        menggambarkan sebuah sikloid.
  \item Rangkuman. Mesin pada Bagian I menghasilkan
        \omterm{pb:b2:curves:1}{astroid}, sebuah parabola semikubik, sebuah \omterm{pb:b2:curves:1}{nefroid} dan sebuah
        sikloid. Untuk masing-masing keempat keluarganya, nyatakanlah dalam satu
        kalimat di mana hipotesis $\Delta \neq 0$ dan
        $E' \neq 0$ berlaku atau gagal, dan peristiwa geometris apa
        (titik runcing, fokus, kemerosotan) yang ditandai tiap kegagalan $E' \neq 0$
        itu. Lalu di mana ekstremum \omterm{thm:b2:curves:frenet2d}{kelengkungannya} mesti muncul pada
        \omterm{pb:b2:curves:1}{selubung} normalnya, dan mengapa?

\end{enumerate}
\end{problem}
